The Discussion interprets the evidence states established in the Results, explains why they
differ across targets, and identifies the limits and next studies implied by unresolved
axes.

\subsection{The contribution is a qualification logic, not another performance benchmark}

We first clarify the study's central contribution and distinguish it from another comparison
of model performance.

The central contribution of this study is to separate signal detection from signal
qualification. Frozen pathology representations made every evaluated association
technically testable, but the scientific status of those associations differed once
transport, confounding, site, representation and endpoint were examined separately. The
result is not a league table of encoders or targets. It is an operational logic for deciding
which claim the available evidence supports, which claim must be restricted to a context,
and which claim remains unsupported in the tested design. The accompanying evidence map
makes that decision explicit and source-linked rather than leaving qualification implicit in
a collection of performance estimates.

\subsection{The three propositions are supported at different evidentiary levels}

We next consider the three propositions separately because the available evidence does not
support each one to the same extent.

The first proposition was supported within the available cross-cohort comparisons: grade
and phenotype, which are directly expressed in routine morphology, showed the clearest
transport. This finding does not establish that morphologic targets are universally easier
than molecular targets, because matching external labels were not equally available.
Rather, it identifies the strongest evidence state attained in this study and demonstrates
what transport evidence looks like when a frozen probe is applied without refitting.

The second proposition was supported by the narrowing of apparently positive signals.
PTEN retained within-cohort support but not an established held-out increment beyond grade.
AR retained a positive pooled direction while site transport and grade-independent
increment remained unresolved. The post-hoc recurrence signal changed with endpoint,
covariate hierarchy and representation. SPOP illustrates a different boundary: the fixed
primary analysis was unsupported and the wider audit was setting-sensitive, but neither
result proves biological absence. Together, these cases show that target morphology,
clinical structure and data provenance determine the scope of a foundation-model claim.

The third proposition was supported by the separation of targets into transportable,
context-sensitive and unsupported evidence states. These states retained information that
a scalar ranking would discard: the same pooled direction could coexist with uncertain
increment, unresolved site transport or endpoint-dependent transfer. The framework is thus
a decision structure for research evidence, not a new significance test or a calibrated
probability that a biomarker is valid.

\subsection{Why pooled performance fails as a scientific claim}

The target-specific results motivate a broader explanation of why a pooled metric cannot, by
itself, support claims about transport, independence, or clinical meaning.

Pooled performance fails when it answers a narrower question than the claim attached to
it. A pooled association can show that information is present in a representation without
showing what generated that information. Grade composition may explain part of a molecular
association; institutional composition may support a pooled effect that is imprecise within
sites; encoder and scale choices may change the recovered signal; and similarly named
clinical outcomes may encode different events. None of these failures makes the original
metric erroneous. They make the broader interpretation underidentified.

This distinction changes how robustness should be reported. Repetition across encoders,
scales, tile budgets and seeds can reveal dependence on analytic choices, but the settings
reuse patients and folds and therefore cannot manufacture independent replication.
Likewise, reproducible code and source-linked outputs establish computational traceability,
not external validity. The practical implication is that a candidate signal should advance
only with the evidence state and unresolved axes reported together. Performance without
that qualification invites the strongest-looking number to stand in for transport,
increment and clinical meaning that were never tested.

\subsection{Limitations and unresolved inferential problems}

The framework also makes visible the evidence that the present datasets and frozen analyses
cannot supply; these limitations bound every qualification decision reported here.

The main limitation is asymmetry of evidence. Grade and phenotype had compatible
cross-cohort labels, whereas molecular targets were available primarily within TCGA-PRAD.
The stronger state attained by morphology-linked targets therefore reflects both their
observed transport and the opportunity to test transport. It cannot establish an inherent
ordering of biological learnability across target classes. The framework also evaluates
only the axes represented in the available data; unmeasured clinical, demographic,
pre-analytic or acquisition factors may remain.

Several uncertainty estimates were inherited from legacy outputs. Some saved intervals did
not retain valid and undefined bootstrap-replicate accounting, and some binary rows did not
retain event counts. Those quantities were not reconstructed. The configuration grid reused
cohorts, folds and patients, so its seed intervals describe tile-sampling variation rather
than patient or population uncertainty. Its cells and paired contrasts are correlated and
cannot be counted as independent confirmations. Consequently, stability across the grid can
bound setting dependence but cannot substitute for a new cohort.

The site analysis cannot separate tissue source from scanner, staining, preparation or case
mix. Site-specific uncertainty therefore remains a transport problem rather than evidence
for a causal acquisition mechanism. Recurrence analyses face a related identification
problem: reconstructed recurrence, progression-free measures and official PFI use different
event definitions and are not interchangeable. The paired analyses used a complete-case
cohort with modest event counts and no multiple imputation, so selection and missingness may
affect the estimated increments.

The recurrence signal was discovered post hoc, and its principal-component analysis (PCA)
sweep was not nested model selection. It therefore remains exploratory even where a point
estimate appears favorable.
More generally, the three evidence states are qualitative research decisions tied to the
frozen analyses and available labels; they are not posterior probabilities, universal
biomarker grades or clinical action thresholds. The study covers prostate pathology, two
frozen encoders and the represented institutions. It does not establish whole-slide
diagnostic performance, treatment utility, population prevalence, causal biology or
generalization to other cancers and deployment environments.

\subsection{Implications and conclusion}

We close by translating the unresolved evidence axes into concrete validation requirements
and a bounded conclusion for pathology foundation-model research.

The unresolved axes specify the next experiments more clearly than a request for more
benchmarking. A transport claim requires an independent cohort with a compatible target
definition and locked preprocessing. A molecular claim requires target-appropriate assays
and patient-level evaluation beyond grade and other clinical structure. A site claim
requires acquisition metadata or a design that can separate site from scanner, stain and
case mix. An outcome claim requires harmonized endpoints, adequate events, model selection
locked before target-cohort evaluation and paired incremental comparison against a clinically meaningful
reference. These are not optional refinements to a pooled metric; they are the evidence
needed to change its qualification state.

The study therefore establishes a concrete rule for pathology foundation-model research:
do not promote a detectable association on pooled performance alone. Report whether it
transports, what clinical information it adds, how it behaves across sites and
representations, which endpoint it predicts, and which dimension remains unresolved. In the
evaluated data, that logic supported grade and phenotype transport, narrowed PTEN and AR,
left SPOP unsupported in the fixed design, and restricted recurrence to an exploratory,
endpoint- and representation-conditioned interpretation. The broader contribution is an
auditable framework and empirical failure map for prioritizing candidate signals without
claiming more than the evidence can support.
